\begin{exercise}{129}
Suppose to be given a commutative diagram
\[
\begin{tikzcd}[column sep=large,row sep=large]
 A_0\arrow[r]\arrow[d,"f_0"] & A_1\arrow[r]\arrow[d,"f_1"] &
 A_2\arrow[r]\arrow[d,"f_2"] & \cdots \\
 B_0\arrow[r] & B_1\arrow[r] & B_2\arrow[r] & \cdots
\end{tikzcd}
\]
in which the maps $f_i$ are Joyal equivalences and all horizontal maps are monomorphisms.
Show that the induced map $f\colon A=\operatorname{colim}A_i\to\operatorname{colim}B_i=B$ is a Joyal equivalence.

\begin{proof}
Let $\mathcal W$ be the class of monomorphisms that are Joyal equivalences;
it is saturated by Exercise~\ref{ex:102}.
The boundary small object argument (Proposition~1.3.9 and
Exercise~\ref{ex:59}) factors any map as a monomorphism followed by a
trivial fibration. The latter is a Joyal equivalence by Proposition~2.2.12.

We construct a factorization of towers $f_\bullet=p_\bullet j_\bullet$,
where $j_n\colon A_n\to C_n$ belongs to $\mathcal W$ and
$p_n\colon C_n\to B_n$ is a trivial fibration.
For $n=0$, use the preceding factorization and two-out-of-three.
Given the construction through $n-1$, form
\[
 P_n=A_n\amalg_{A_{n-1}}C_{n-1}.
\]
The map $d_n\colon A_n\to P_n$ is a pushout of $j_{n-1}$, hence lies
in $\mathcal W$. The maps $f_n$ and
$C_{n-1}\xrightarrow{p_{n-1}}B_{n-1}\to B_n$ induce
$v_n\colon P_n\to B_n$ with $v_nd_n=f_n$, so $v_n$ is a Joyal equivalence.
Factor $v_n=p_n\lambda_n$ by the boundary small object argument.
Then $\lambda_n\colon P_n\to C_n$ lies in $\mathcal W$ by two-out-of-three.
Set $j_n=\lambda_nd_n$ and use
$C_{n-1}\to P_n\xrightarrow{\lambda_n}C_n$ as transition map.
The pushout identities make $j_\bullet$ and $p_\bullet$ maps of towers.

Put $C=\operatorname{colim}_nC_n$ and $D_n=A\amalg_{A_n}C_n$.
The map $A\to D_0$ is a pushout of $j_0$.
For $n\ge1$, the canonical maps $P_n\to D_{n-1}$ give pushout identities
\[
 D_{n-1}\amalg_{P_n}C_n\cong A\amalg_{A_n}C_n=D_n.
\]
Thus each $D_{n-1}\to D_n$ is a pushout of $\lambda_n$.
Taking colimits of the defining pushouts gives
$\operatorname{colim}_nD_n\cong C$, since $\operatorname{colim}_nA_n=A$.
Saturation of $\mathcal W$ now shows that $A\to C$ belongs to $\mathcal W$.

Finally, $C\to B$ is a trivial fibration.
Indeed, in a boundary lifting problem, the two horizontal maps factor
through stages by Exercise~\ref{ex:55}, since $\partial\Delta^m$ and
$\Delta^m$ have finitely many nondegenerate simplices, also for $m=0$.
The same exercise makes the square commute at a common later stage.
There $p_n$ supplies a filler, which passes to the colimit.
Hence $f\colon A\to C\to B$ is a composite of Joyal equivalences.
\end{proof}
\end{exercise}
